\subsection*{Espaços Vetoriais}

\begin{frame}[label=vetores]{Espaços Vetoriais Reais}

Um {\color{blue}espaço vetorial real} $V$ é um conjunto, cujos elementos são chamados {\color{blue} vetores}, no qual estão definidas duas operações:

\begin{itemize}
\item A {\color{blue} adição}, que a cada par de vetores $\vec{u}$,$\vec{v}\in V$ faz corresponder um novo vetor $\vec{u}+\vec{v}$, chamado {\color{blue}vetor soma} de $\vec{u}$ e $\vec{v}$.

\item A {\color{blue} multiplicação por escalar}, que a cada número $\alpha\in \R$ e a cada vetor $\vec{v}\in V$ faz corresponder um vetor $\alpha \vec{v}$, chamado {\color{blue}produto de $\alpha$ por $\vec{v}$}.
\end{itemize}
Além disso, essas operações devem satisfazer, para quaisquer $\vec{u},\vec{v},\vec{w}\in V$ e $\alpha,\beta \in \R$, as seguintes condições:
%\begin{enumerate}
%\item $\vec{u}+\vec{v}=\vec{v}+\vec{u}$ (comutatividade)
%\item $(\vec{u}+\vec{v})+\vec{w}=\vec{u}+(\vec{v}+\vec{w})$ (associatividade)
%\item Existe um vetor $\vec{0}\in V$, chamado {\color{blue}vetor nulo}, tal que
%\[\vec{u}+\vec{0}=\vec{u} \text{ para todo } \vec{u}\in V.\]
%\end{enumerate}
\end{frame}



\begin{frame}[label=vetores]{Espaços Vetoriais Reais}

\begin{enumerate}

\item $\vec{u}+\vec{v}\in V$. (fechamento em relação à adição)
\item $\vec{u}+\vec{v}=\vec{v}+\vec{u}$ (comutatividade)
\item $(\vec{u}+\vec{v})+\vec{w}=\vec{u}+(\vec{v}+\vec{w})$ (associatividade)
\item Existe um vetor $\vec{0}\in V$, chamado {\color{blue}vetor nulo}, tal que
\[\vec{u}+\vec{0}=\vec{u} \text{ para todo } \vec{u}\in V.\]

\item Para cada $\vec{v}\in V$, existe um vetor $\vec{-v}\in E$, chamado {\color{blue}inverso aditivo}, ou {\color{blue}simétrico} de $\vec{v}$, tal que
\[-\vec{v}+\vec{v}=\vec{0}.\]

\item $\alpha \vec{v}\in V$. (fechamento em relação à multiplicação por escalar)

\item $(\alpha+\beta)\vec{v}=\alpha \vec{v}+\beta \vec{v}$ e $\alpha(\vec{u}+\vec{v})=\alpha\vec{v}+\alpha\vec{v}$. (distributividade)

\item $1\cdot \vec{v}=\vec{v}$.
\end{enumerate}
\end{frame}



\subsection*{O Espaço Vetorial \(\R^n\)} \label{evreal}
\begin{frame}[label=vetores]{O Espaço Vetorial \(\R^n\)}

O conjunto da matrizes $\mathbb{M}_{m\times n}(\R)$,
com as operações de soma e multiplicação por escalar, formam um {\color{blue} espaço vetorial real}.
\medskip

Em particular, o conjunto das \dt{matrizes coluna} (ou matrizes linhas) será denotado por $\R^n$,
e a seus elementos chamaremos \dt{vetores}. Neste caso, usaremos a seguinte notação \(\vec{v}\in\R^n \), o que significa 
\[
\vec{v}=\begin{bmatrix}
v_1\\ v_2\\ \vdots\\ v_n
\end{bmatrix} \text{ ou } \vec{v}= [x_1,x_2,\ldots,x_n],
\]
tais que $v_i\in \R$ para todo $i=1,2,\ldots,n$.
\end{frame}

\begin{frame}[label=vetores]
Um nutricionista forneceu uma tabela que especifica a {\color{blue} quantidade mínima de cada tipo de vitamina que deve ser ingerida diariamente}.

\begin{center}

\begin{tabu}{|l|c|c|c|c|}
\hline
Tipo de vitamina & A & B & C & E\\ \hline
 Quantidade  mínima (mg) & 3 & 1.1 & 60 & 11\\ \hline
\end{tabu} 
\end{center}

Ele também forneceu uma tabela com a quantidade (em mg) de vitaminas em cada 100 gramas de 4 tipos de alimentos diferentes,
\begin{center}
\begin{tabu}{|c|c|c|c|c|}
\hline
& \multicolumn{4}{c|}{Alimento} \\ \hline
Vitamina & 1 & 2 & 3 & 4  \\ \hline
 \rowfont{\color{blue}} A &  0.140 & 0.580  & 0.150 & 0 \\ \hline
 \rowfont{\color{blue}} B &  0.08  & 0      & 1.30  & 0.08  \\ \hline
 \rowfont{\color{blue}} C &  0     & 0      & 26    & 38  \\ \hline
 \rowfont{\color{blue}} E & 1.60  & 0      & 6.90    & 0.2 \\ \hline
\end{tabu}
\end{center}

 Cada alimento é uma {\color{red}variável}, ou seja, 
\begin{center}
 $x_i$ representa a quantidade (em ``pacotes'' de 100g) do alimento $i$ que deve ser consumido diariamente, com $i=1,\ldots,4$. 
 \end{center}\bigskip

\end{frame}

\begin{frame}[label=vetores]
Com isso, se quisermos saber qual a quantidade de cada alimentos devemos consumir, para se ter a quantidade mínima de cada vitamina, devemos resolver o sistema:
\[
\begin{cases}
{\color{blue}0.14 } {\color{red}x_1}+ {\color{blue}0.58}{\color{red}x_2}+ {\color{blue}0.15}{\color{red}x_3}= 3,\\
{\color{blue}0.08}{\color{red}x_1}+{\color{blue}1.3}{\color{red}x_3}+{\color{blue}0.08}{\color{red}x_4}= 1.1,\\
{\color{blue}26}{\color{red}x_3}+{\color{blue}38}{\color{red}x_4}=  60,\\
{\color{blue}1.6}{\color{red}x_1}+{\color{blue}6.9}{\color{red}x_3}+{\color{blue}0.2}{\color{red}x_4}=  11,
\end{cases} 
\]
Usando a notação vetorial, temos:

\[
{\color{blue}
\begin{bmatrix}
0.14 \\ 0.08 \\ 0 \\ 1.60
\end{bmatrix}} {\color{red}x_1}+
{\color{blue}
\begin{bmatrix}
0.58 \\ 0 \\ 0 \\ 0
\end{bmatrix}} {\color{red}x_2}+
{\color{blue}
\begin{bmatrix}
0.150 \\ 1.30 \\ 26 \\ 6.90
\end{bmatrix}} {\color{red}x_3}+
{\color{blue}
\begin{bmatrix}
0\\ 0.08 \\ 38 \\ 0.2
\end{bmatrix}} {\color{red}x_4}=
\begin{bmatrix}
3 \\ 1.1\\ 60 \\ 11
\end{bmatrix}
\]
\end{frame}

\begin{frame}[label=vetores,fragile=singleslide]
\begin{pycode} 
import sympy as sp

x1,x2,x3,x4=sp.symbols('x_1 x_2 x_3 x_4',real=True)
X=sp.Matrix([x1,x2,x3,x4])
def pt(Y):
 return((Y[0],Y[1]))

v1=sp.Matrix([0.14,0.08,0,1.6])
v2=sp.Matrix([0.58,0,0,0])
v3=sp.Matrix([0.15,1.3,26,6.9])
v4=sp.Matrix([0,0.08,38,0.2])
B=sp.Matrix([3,1.1,60,11])

A=sp.Matrix.hstack(v1,v2,v3,v4)

AX=A*X
sol=sp.solve(sp.Eq(AX,B),(x1,x2,x3,x4))
\end{pycode} 
Aprenderemos mais à frente que este sistema tem a seguinte solução
\[x_1 \approx\pyl{sp.N(sol[x1],3)},\ x_2 \approx\pyl{sp.N(sol[x2],3)},\ x_3 \approx\pyl{sp.N(sol[x3],3)}\ \text{ e }\ x_4 \approx\pyl{sp.N(sol[x4],3)}.\]
Isto é, para se {\color{blue}ingerir a quantidade mínima de cada tipo de vitamina} devemos consumir:
\begin{center}

$\pyl{sp.N(sol[x1]*100,3)}$ gramas do alimento 1\\
$\pyl{sp.N(sol[x2]*100,3)}$ gramas do alimento 2\\
$\pyl{sp.N(sol[x3]*100,3)}$ gramas do alimento 3\\
$\pyl{sp.N(sol[x4]*100,3)}$ gramas do alimento 4
\end{center}
\end{frame}


